\subsection{Effect of replay subtraction}\label{sec:update-ablation}
To distinguish update direction from magnitude, we compare the paired update with direct physical fine-tuning and a norm-matched direct update. The latter is $\theta_{\rm control}=\theta_{\rm base}+\alpha(\theta_{\rm physical}-\theta_{\rm base})$, where $\alpha$ is the ratio of paired to direct Euclidean parameter-update norms, excluding buffers. The coefficients, computed before evaluation, are 0.2536 and 0.2564.

\begin{table}[t]
\centering
\caption{Update controls on 10 compositions, with 1,280 attempts per model. eSEN energy changes are relative to the base model on jointly graph-valid output pairs.}
\small
\begin{tabular}{lrr}
\toprule
Model & Graph validity (\%) & Energy change (eV/atom)\\
\midrule
Base &62.11&0\\
Direct physical update &54.77&$+0.11080$\\
Same-norm direct update &61.64&$+0.01893$\\
Paired update &62.03&$-0.01554$\\
\bottomrule
\end{tabular}\label{tab:update-ablation}
\end{table}

The norm-matched direct update largely preserves graph validity but raises energy, whereas the paired update lowers it. Paired minus norm-matched direct gives energy change $-0.03232$ eV/atom ([-0.03981,-0.02485]) and force-RMS change $-0.4170$ eV/\AA\ ([-0.5529,-0.2859]). Both runs and the all-output comparison favor the paired direction. Graph validity differs by $+0.39$ points ([-1.95,2.73]). The physical benefit therefore depends on update direction, not merely a smaller parameter change.
